\section{Derivations}
\label{app:derivations}

This appendix collects the sourced values and working behind choices the text states in words. The input table holds the ranges (\cref{tab:inputs}), and the published code is the full record of every number.

\subsection{Calm-water calibration}
\label{app:calm_calibration}

MAN rates its bulk carriers at 85\,\% of the rating up to \num{84000}~t and at 75\,\% for the Capesize \citep{man2022bulk}, and the MR tanker \citep{man2019mr} and the container ships \citep{man2024container} at 90\,\%, all with $s=15\,\%$. For the Aframax tanker and the VLCC, whose sources state no margins, we take the traditional 10\,\% engine margin and 15\,\% sea margin \citep{man2011}. The car carrier's source design makes \SI{15.0}{kn} at 43\,\% of its rating with a 15\,\% sea margin \citep{keh_8600ceu_df_pctc}, 39\,\% at its \SI{14.5}{kn} by the cube law.

Before calibrating we set the propeller diameter by type and the longitudinal centre of buoyancy, the centre of the underwater volume, by Schneekluth's regressions: $\mathrm{lcb}=-13.5+19.4\,C_P$\,\% of the length forward of midships for tankers and bulk carriers and $8.80-38.9\,\mathrm{Fn}$\,\% for the others \citep{ventura_estimation}, with the Froude number $\mathrm{Fn}=V/\sqrt{gL}$. The larger propeller removes 8 to 23 points of the excess and leaves 1 to 29\,\% on ten classes. The rest is the method itself on full hulls: a form factor of 1.31 to 1.39 on the bulk carriers and tankers, and a correlation allowance of 11 to 17\,\% of the resistance. On the Ultramax and Kamsarmax, whose block coefficients of 0.86 and 0.88 exceed the 0.75 to 0.85 that \citet{man2011} lists for bulk carriers, the forward centre of buoyancy also drives the wave-making coefficient to 8 and $10\times10^{-4}$, and the excess to 61 and 84\,\%.

\subsection{Cargo units}
\label{app:cargo_units}

Containers: 18.8\,\% of the TEU handled in EU main ports were empty in 2024 \citep{eurostat_mar_mg_am_cvh}. A loaded TEU carries \SI{10.2}{t} of goods, the mean over 2015--2024 of the gross weight of goods in large containers handled in EU main ports \citep{eurostat_mar_mg_am_cwhc} per loaded TEU \citep{eurostat_mar_mg_am_cvh}. This weight excludes the box \citep{eurostat_maritime_manual}, so we add \SI{1.85}{t} of tare, half a \SI{40}{ft} box \citep{hapaglloyd_container_spec,eurostat_maritime_manual}, in \SI{38.3}{m^3}, the external box of ISO 668 that \citet{mao2020icct} count a fuel tank against.

Car decks: a CEU carries \SI{1.3}{t}, since a PCTC of \num{8600}~CEU is optimised for \SIrange{8000}{15000}{t} of payload \citep{keh_8600ceu_df_pctc}. As a check, new cars in Europe weigh \SI{1.56}{t} in running order \citep{icct2025_eea_brief} and a car takes about 1.3~CEU \citep{portantwerpbruges2024_ceu}, or \SI{1.2}{t} per CEU. A CEU takes \SI{8.4}{m^2} of deck \citep{keh_8600ceu_df_pctc} at a deck pitch of \SI{2.9}{m}, \SI{24}{m^3}.

\subsection{Published kite aspect ratios and areas}
\label{app:kite_aspect}

\citet{paulig2013} give the SkySails ship kite a flat size of \SI{225}{m^2} and a projected aspect ratio below 2, chosen for its tolerance to changes in angle of attack and its stability at launch and landing, at a glide ratio below 5. They report that the gain in availability outweighed the loss in aerodynamics, without numbers. A ram-air wing from SkySails Power, \SI{125}{m^2} projected of \SI{160}{m^2} total with a projected span of \SI{15.2}{m}, has a projected aspect ratio of 1.86 \citep{folkersma2020}. That wing is not stated to be a ship kite, and \citet{paulig2013} give about \SI{14}{m} for the span of their \SI{160}{m^2} kite. \citet{penloup2021} describe the SkySails ship kites as of moderate and very low aspect ratio.

SkySails flew its kite of \SI{160}{m^2} flat, \SI{125}{m^2} projected, on \SI{300}{m} \citep[Fig.~18]{erhard2013} and \citet{paulig2013} put the operating point of their \SI{320}{m^2} kite, about \SI{250}{m^2} projected, at about \SI{360}{m}. An aspect ratio of 3 at those areas makes these 15.5 and 13.1 spans. On its own \SI{15.2}{m} span the first would be 19.7, about the upper bound of our range. The SkySails tether weighs \SI{0.9}{kg/m} for \SI{320}{kN} \citep{paulig2013}, against \SI{0.66}{kg/m} for plain rope here.

\subsection{Tether life}
\label{app:tether_life}

\Cref{tab:inputs} gives the two creep laws and their bounds: the power law of \citet{humeau2018}, anchored at the mooring-line creep life of \citet{dsm_sk78_mooring}, and the diagram of \citet{bosman2013}. Load peaks count in the sum, since the tension fatigue of this rope is creep under the varying load \citep{humeau2018}. They are the swing of each figure eight, gusts at our estimated turbulence intensity of 6\,\% (4 to 10\,\%), and the bow's motion, which the stretch of the rope turns into a tension change of at most 5\,\%. A thicker tether at \SI{0.3}{GPa} would last about \num{50000} hours by the same laws, and a thinner one at \SI{0.48}{GPa} about \num{7300}.

Bending under load would bind much sooner. Pumped over a drum at this design factor a tether fails after 45 days, about \num{1100} flight hours \citep{bosman2013}, and a rope bent and straightened at 20 diameters under a third of its break load lasts about \num{2500} bends \citep{davies2015}. Tug lines led through a smooth staple keep 70\,\% of their strength after 300 jobs and 42\,\% after \num{1250}, from sand and salt crystals worked into the rope and bending at the points of contact \citep{mccorkle2004}. Sea water does not weaken this fibre \citep{dsm_sk78_mooring}. HMPE ship mooring lines, also held at a fixed length at sea and led through a fairlead, are designed for \num{8000} mooring hours to 75\,\% of their strength at up to 22\,\% of break load and retired for wear in port rather than creep \citep{wilhelmsen2022}, so their life does not bound ours. The fixed tether section of the Ampyx Power design, replaced after 7 years \citep{vanhagen2023}, is an assumption of that study with no stated basis.

\subsection{Tether bending at launch and landing}
\label{app:tether_bending}

We carry the law of \citet{davies2015}, measured on a \SI{19}{mm} rope over a sheave of 20 rope diameters, down to the reeling tension of a tenth of the force limit, about 3\,\% of break load, and up to a drum of 40 diameters with the factors of \citet{bosman2013}. As a power law it gives \num{840000} bends, with the slope of Bosman's tests in log life over tension \num{160000}, and we take $N_b=\num{400000}$. The low bound, \num{5500}, reels at 10\,\% of break load over 20 diameters, scaled to a \SI{70}{mm} tether by the size effect between Bosman's 8 and \SI{20}{mm} ropes and cut 30\,\% for reverse bending. The high bound, \num{6} million, reels at 1.7\,\% over 100 diameters. Bending and creep damage add by Miner's rule.

\subsection{Canopy mass against the SkySails kite}
\label{app:canopy_mass}

With the cloth of slightly above \SI{100}{g/m^2} of the SkySails kites, our build-up would give \SI{0.40}{kg/m^2}, against the \SI{0.68}{kg/m^2} projected of their \SI{225}{m^2} flat kite of \SI{120}{kg} with its lines \citep{paulig2013}. Beyond about \SI{66}{kg} of cloth and \SI{18}{kg} of lines, both our estimates, that kite leaves \SI{0.19}{kg/m^2}. It was built for about \SI{0.9}{kN/m^2} projected, so in proportion to the limit its \SI{0.19}{kg/m^2} would be about \SI{0.11}{kg/m^2} at the base limit, and the high end of this input therefore overstates what it reports. That high end turns on our two estimates of its cloth and lines. With the model's own bridle for that kite, about 2~kg, and its area at the 1.28 flat per projected of \citet{folkersma2020}, the same weight leaves about \SI{0.30}{kg/m^2} and a canopy of \SI{1.00}{kg/m^2}, against \SI{0.89}{kg/m^2} at the high end. Seams, tapes and reinforcements stay in the making cost, about USD~3.5 of USD~42 per square metre, so a kite built for twice the limit would pay about 5\,\% more for its canopy than we count. A kite built for \SI{1}{kN/m^2} carries \SI{241}{g/m^2} of cloth and one built for \SI{0.3}{kN/m^2} \SI{174}{g/m^2}.

\subsection{Canopy prices}
\label{app:canopy_price}

\begin{table}[tbp]
\centering
\footnotesize
\caption{Canopy price today and at the goal, USD per projected square metre unless stated. Sourced inputs are in \cref{tab:inputs}.}
\label{tab:canopy_price}
\begin{tabular}{@{}lrr@{}}
\toprule
& Today & Goal \\
\midrule
Cloth, USD per sewn m$^2$ of cloth & 17 & 9.4 \\
Cloth per projected m$^2$ & 60 & 33 \\
Sewing, hours per flat m$^2$ & 1.5 & 1.0 \\
Making cost & 57 & 42 \\
Cost to make & 117 & 75 \\
Price at a market factor of 2 & 234 & 150 \\
\bottomrule
\end{tabular}
\end{table}

Today's hours are the low end of what a maker's 50 hours per paraglider gives \citep{skycz2017}. Polyester sailcloth of 170 to \SI{215}{g/m^2} costs USD~11 to 19 per square metre at retail without VAT \citep{extremtextil_prices_2026,sailrite_dacron_prices_2026}, and USD~9.4 stays above the airbag-grade weaving and solution-dyeing premium for which we have a price. The goal's 1.0 hour is inferred from the G-12 cargo parachute in a year of stock-building, when a USD~19.2~million order bought \num{5566} complete assemblies \citep{bis2004,mills_g12_web} (\cref{tab:inputs}). At the goal a canopy costs USD~43 per square metre of finished cloth, above the roughly USD~12 of a complete nylon G-12 parachute in 2002 \citep{bis2004,mills_g12_web} and below the EUR~49 (USD~57) per square metre from which a sail loft offers a crosscut Dacron mainsail \citep{horizonsails_web}.

\subsection{Seams and cloth of a ship kite}
\label{app:seam_cloth}

The figures below are our own estimates on read inputs. Rib cloth follows the proportion of a cell and seam length its width, so a ship kite cuts seams far more than cloth.

\begin{table}[tbp]
\centering
\footnotesize
\caption{Cell layout of a paraglider and of ship kites. \tbd{Paraglider data: cite Ozone Rush 5.} The SkySails cell count is from \citet{paulig2013}, the widths by our reading.}
\label{tab:seam_cloth}
\begin{tabular}{@{}lrrr@{}}
\toprule
& Paraglider & SkySails \SI{225}{m^2} flat & This kite \SI{800}{m^2} \\
\midrule
Flat aspect ratio & 5.55 & -- & 3.76 \\
Cells & 57 & 22 & 31 \\
Cell width (m) & 0.21 & 1.1 & 1.9 \\
Cell length over width & 10.3 & 8.1 & 8.1 \\
Seam per flat m$^2$ (m) & 14 & -- & 2.5 to 4.5 \\
\bottomrule
\end{tabular}
\end{table}

No source says what share of the hours follows seam length. At shares of 20, 50 and 80\,\% the 1.5 hours become about 1.3, 0.9 and 0.6 hours, before heavier cloth, load seams and the handling of large panels, none of them measured. \citet{dunker2013} supports the seam basis without giving a rate, as shorter curved skin seams make sewing quicker. He recommends about \SI{115}{kg} at most per subassembly under a sewing machine, so our \SI{600}{kg} canopy of \SI{800}{m^2} takes at least six modules, with joins and handling for which he gives no hours. We do not read across from the \SI{12}{m^2} leading-edge inflatable kite at 3.7 hours per square metre: its 44 hours against the paraglider's 50 on twice the area look like hours per kite, from tubes and bladders a ram-air kite does not have.

For the cloth, two skins take about \SI{2.05}{m^2} per flat square metre and ribs and diagonals scale with the proportion, fitted to the paraglider. This gives \SI{3.3}{m^2} per projected square metre at 31 cells (2.9 to 3.9), 3.0 at 22 cells on our wider planform, and a floor of 2.4 at no ribs. Cell width is typically below section thickness \citep{dunker2013}, about 2.2 to \SI{2.5}{m} here. The \SI{1.9}{m} of 31 cells meets that and the \SI{2.7}{m} of 22 cells does not, so 3.3 is the more likely. The cost of fewer cells that \citet{paulig2013} report, stronger cloth almost cancelling the saved mass, does not bind here, where life and not strength sets the cloth weight. The variants take the following off the goal price and the canopy mass. Hours matter more for the price and only the ratio touches the mass. The spread of the savings in \cref{tab:missed_goals} suggests 0.2 to 0.4 points of saving for USD~9 and 0.4 to 1.0 for USD~24.

\begin{table}[tbp]
\centering
\footnotesize
\caption{Effect of a lower cloth ratio or fewer sewing hours on the goal price and the canopy mass.}
\label{tab:canopy_variants}
\begin{tabular}{@{}lrr@{}}
\toprule
Variant & Price (USD/m$^2$) & Canopy mass (kg/m$^2$) \\
\midrule
Base & 150 & 0.75 \\
Cloth ratio 3.0 & $-9$ & 0.65 \\
Cloth ratio 2.4, no ribs & $-21$ & 0.53 \\
Sewing 0.8 h & $-12$ & 0.75 \\
Sewing 0.6 h & $-24$ & 0.75 \\
\bottomrule
\end{tabular}
\end{table}

\subsection{Learning curve check of the canopy price}
\label{app:learning_curve}

Under Wright's law the time per unit falls by a fixed ratio $p$ with each doubling of cumulative output. For sewing $p=0.93$ for the sewing section of a shirt plant over 15 shirts \tbd{cite badri2016} and 0.77 to 0.83 for three experienced machinists on a new task over 20 cycles \tbd{cite yepez2025}, computed from the fitted exponents because the percentages printed there are not per doubling. We take 0.83 and 0.93 as the ends, leaving out the fastest machinist since a check should not rest on one person. Both ends lie inside the 0.75 to 0.95 called typical of industry \tbd{cite tamas2020}. Learning also stops somewhere, more readily the more machine-paced the work \tbd{cite tamas2020}, and we test both no floor and 0.6 hours, the bottom of the G-12 range. The curve starts at a volume, since the 1.5 hours come from a series in production: 0.1 and 1 million flat square metres, \num{4000} and \num{40000} gliders of \SI{25}{m^2}, a new plant inheriting part of that experience and the whole history of an established maker.

At \num{1000} canopies of \SI{800}{m^2} this gives 0.80 to 1.44 hours per flat square metre and at \num{10000} canopies 0.44 to 1.30 hours. With cloth at the goal and USD~6 for each 0.1 hour, the price is USD~138 to 177 and USD~116 to 169. If seams are a fifth of the hours, the curve starts at about 1.25 hours and needs 140 to 800 canopies from the low start. If they are half or more, the goal needs no learning. Polyester fibre has $p=0.72$ in a dataset whose chemicals come largely from a 1972 compilation \tbd{cite lafond2018}, but it counts world output, 78 million tonnes in 2024 \tbd{cite Textile Exchange 2025}, against a fleet's 280 to 560 tonnes a year, and yarn is USD~0.4 of the USD~8.2 before offcuts.

\subsection{Data behind the canopy life}
\label{app:canopy_life_inputs}

\paragraph{UV} A least-squares fit of the strength loss of polyester without a UV blocker over \SIrange{40}{1500}{g/m^2} gives $a=0.40\pm0.12$ and $D_\mathrm{ref}=\SI{63}{MJ/m^2}$, from outdoor tests with a measured dose between 295 and \SI{385}{nm} \citep{lodi2008,stevens2018}, a roof test \citep{hunt2022}, and lamp tests \citep{ceylan2023,asadi2021,uzumcu2024,dekanic2023} times 3, the gap on the one cloth tested both ways \citep{stevens2018}. The field values alone give 0.42 and \SI{72}{MJ/m^2}. The bounds are 0.1 and 1 on $a$ and 35 and \SI{160}{MJ/m^2} on $D_\mathrm{ref}$. The yearly mean dose in Florida is \SI{320}{MJ/m^2} between 295 and \SI{385}{nm} spread over every hour \citep{qlab2026}.

\paragraph{Wear} The statements span 0.04 to \SI{0.26}{g/m^2/h}, and the sail's shape life is \numrange{1700}{2700} hours \citep{flynn2015}. The bounds on $r$ are 0.052 and 0.26. Cloth and sail both work at 0.4 to 3\,\% of their breaking strength. Rope yarns cycled near that load lost 0.4 to 9.2\,\% in $10^6$ cycles \citep{tti2006}, and woven polyester kept its strength over \num{12000} cycles to half or three-quarters of its breaking strain \citep{nosraty2013}. Flutter could supply the rate at about 1\,\% of flight time in the one test with a stated duration \citep{doyle2011}, and folds at hundreds of crumples an hour \citep{maslanka2023}. The sailmakers' line that 30 minutes of flogging can equal 50 hours of sailing \citep{allatsea2007} has no test behind it and is not used.

\paragraph{Exponent $b$} The bounds are 0 and 1.5. One maker's lives for its balloon cloths imply 1.5 \citep{kubicek2026}, one life for every woven polyester sail implies about 0 \citep{flynn2015}, and the one measurement, on nylon cloths of 34 to \SI{42}{g/m^2}, finds no order with weight \citep{maslanka2023}. With the life at \SI{200}{g/m^2} held, these bounds put the life at \SI{100}{g/m^2} at 32 to 73\,\% of it.

\paragraph{Tension over cloth weight} No data set varies cloth weight and tension with strength loss measured at a canopy's tension, since cloth is chosen for its load. Wings at $T/w$ of 1 to 3 live from under 130 flying hours as skydiving mains, at \num{1600} jumps \citep{szafran2018} of 5 minutes \citep{szafran2019}, to the unlimited life stated for fully coated aircraft polyester \citep{faa2001}. Folding, opening shock, flutter and UV sort these lives and $T/w$ does not.

\paragraph{End of life} One tandem paraglider fell from about 400 to \SI{29}{s} in the porosity test by 800 hours, against a limit of \SI{25}{s} \citep{supair2016}. Other makers' limits are \SI{18}{s} \citep{dudek2026} and \SI{10}{s} \citep{bgd2026}. A sail keeps its shape for half to two-thirds of its structural life \citep{flynn2015}, and a paraglider maker limits a tear to \SI{5}{mm} at a load of \SI{600}{g} \citep{bgd2026}. The margin of \cref{eq:uv_life} flies \SI{200}{g/m^2} down to 4\,\% of new strength. Aircraft fabric is retired at 70\,\% of the new strength required, by a rule tabulated for cotton and linen \citep{faa2001}, and balloon cloth at a test value of 46 and 27\,\% of new, at most \citep{lindstrand_mm}. The retirement bounds of \cref{sec:det_canopy_life} rescale $r$ to keep the observed wear life at \SI{100}{g/m^2}. Thread is not shown to fail before cloth on the one vendor curve available \citep{shippee2016}, though a sailmaker reports stitching as the weakest part of a sail \citep{gifford2014}. The launch count of \num{2500} takes 1.5 hours a flight, the ratio of the more than \num{1800} hours of air time to the more than \num{1200} automated flights SkySails reported over all its kites \citep{skysails2018}.

\subsection{Range of the canopy life}
\label{app:canopy_life_range}

\begin{table}[tbp]
\centering
\footnotesize
\caption{Canopy life at \SI{200}{g/m^2} with one input at each bound, in flight hours.}
\label{tab:canopy_life_range}
\begin{tabular}{@{}lrr@{}}
\toprule
Input changed alone & Low & High \\
\midrule
Wear rate $r$ & \num{2200} & \num{6900} \\
Reference dose $D_\mathrm{ref}$ & \num{3100} & \num{4400} \\
UV exponent $a$ & \num{3600} & \num{4200} \\
Dose rate $\dot D$ & \num{3200} & \num{4200} \\
Cloth strength $\sigma_w$ & \num{3300} & \num{4100} \\
Kept strength of 8 and 3 times the tension & \num{3300} & \num{4400} \\
Wear exponent $b$ of 0 and 1.5 & \num{2200} & \num{4900} \\
Retirement at 70 and 27\,\% of new strength & \num{1200} & \num{2400} \\
\bottomrule
\end{tabular}
\end{table}

A fixed wear life of \num{2000} hours, about the shape life of a sail, gives \num{1800}. With every input at its bound together the life is \num{700} and \num{18000} hours, but no data say how the inputs vary together.

\subsection{Cloth tension data}
\label{app:cloth_tension}

The \SI{3200}{m^2} kite in \SI{12}{m/s} of apparent wind, \SI{90}{\degree} off downwind, flies a cycle of \SI{82}{s}. Its tether force runs from 1.12 to \SI{1.60}{MN}, with a mean of \SI{1.52}{MN}, and stays at the force limit for 68\,\% of the cycle. The cloth carries \SI{237}{N/m} on average and new cloth holds \SI{34}{kN/m}. The radius of \SI{0.41}{m} averages over the chord of a section 15\,\% thick, our assumption, on the SkySails kite of \SI{125}{m^2} projected, with 38 cells over \SI{15.2}{m} of projected span and a largest chord of \SI{9.26}{m} \citep{folkersma2020}, about \SI{0.43}{m} wide once flattened by our estimate.

\subsection{VTOL kite data}
\label{app:vtol_data}

\paragraph{Launch alternatives} A rotating arm needs a radius of at least \SI{30}{m} even at \SI{15}{kg/m^2} and a winch launch a run of about \SI{130}{m} \citep{fagiano2017takeoff}, while these wings carry 34 to \SI{80}{kg/m^2}. A catapult or rail with a capture system must also line up with the relative wind \citep{joshi2025}. The one offshore flight was from a moored spar buoy in \SI{10}{m/s} and \SI{1}{m} waves \citep{echeverri2020part1}.

\paragraph{Mass scaling} Wing~7 has \SI{60}{kg} on \SI{3.96}{m^2} and \SI{8}{m} of span \citep{vanderlind2013}. To MX2 it gives 2.91 by span and 2.63 by area, which differ by the aspect ratio of 16 against 12.5. To the built M600, of 1692 to \SI{1731}{kg}, it gives 2.87 to 3.18. For wings launched without propellers \citet{sommerfeld2022scaling} sweep 2.7 to 3.3 as an assumption and \citet{vanhagen2023} use 2.7 to 3.0. The designs of \citet{joshi2025} run from \SI{700}{kg} at \SI{20}{m^2} to \SI{10663}{kg} at \SI{160}{m^2}, with the remark that a VTOL wing will scale differently because of its motors. From the same anchor the \SI{150}{m^2} design of \SI{6885}{kg} of \citet{eijkelhof2020reference} gives 2.57. The tether holds \SI{5.56}{kN/m^2} of wing at a working stress of \SI{0.44}{GPa} and is 11.5 spans long, with half its weight hanging at the kite. On a fixed length a tether's drag coefficient would fall as the wing grows \citep{ruiterkamp2013}.

\paragraph{Propeller drag} Generating costs a drag coefficient of 0.09 to 0.18 on the area of MX2 in the propeller table of the release, windmilling 0.019 to 0.022 at 40 to \SI{50}{m/s}, and a propeller held slow 0.10 to 0.13. With one propeller out in hover each of the others needs 107~kW and \SI{637}{N.m}, against 112~kW and \SI{675}{N.m} in generation.

\paragraph{Loops} The much lighter PowerPlane of Ampyx, \SI{28}{kg} on \SI{3}{m^2}, did fly figure eights \citep{ruiterkamp2013}. The wings that flew used 6.25 spans on Wing~7 \citep{vanderlind2013} and 5.65 on the M600, on tethers of 16.5 to 20 spans. \citet{ruiterkamp2013} expect a massless balance to overstate the power of a wing that keeps accelerating, while \citet{vanderlind2013} finds an analysis of this kind within about \SI{1}{m/s} of the flown power curve of Wing~7, at \SI{15}{kg/m^2}, against the \SI{34}{kg/m^2} of MX2. Wing~7's control held a lowest speed of \SI{20}{m/s}, 1.8 times its stall speed. On the published inputs of Wing~7 this model flies from \SI{6.5}{m/s} with the stall speed and \SI{7.5}{m/s} with that margin, while its published curves pass zero power at about \SI{5.3}{m/s}, read from the figure, and flights lie about \SI{1}{m/s} either side.

\paragraph{Hover} Momentum theory uses propellers sweeping 0.61 of the wing area, a thrust of 1.25 times the weight on six of eight propellers \citep{echeverri2020part1} and an electric to hover efficiency of 0.62 \citep{echeverri2020part1,makani2020code}. It gives \SI{13.3}{MW} for \SI{300}{m^2}.

\paragraph{Prices} The structure costs USD~290 per kilogram and USD~232 per square metre, with USD~\num{174000} of avionics \citep{joshi2025}. The drive on board costs USD~0.10 per watt at 1.33 times the hover power, two thirds of the estimate of \citet{bauer2018}, whose last third is the electronics on the ground, here the ship's supply for the hover peak at our estimate of USD~0.15 per watt. The electrified tether costs USD~186 per kilogram, twice the rope of the soft kite, and a perch USD~\num{100000}, our estimate. Upkeep is USD~0.61 per square metre and flight hour, the most that Makani allows a kite system for all its maintenance \citep{echeverri2020part1}, so the airframe pays no maintenance on its price besides.

\subsection{Convergence of the routing search}
\label{app:routing_convergence}

Both searches are on the 12 departures of the results for the \SI{140}{m} e-ammonia ship and the Kamsarmax on VLSFO, each bare and with each rig, on the routes to Guadeloupe, New York and Tubarão. The backward estimate runs at times \SI{48}{h} apart, on every second commanded speed, with the weather at most \SI{6}{h} apart along an edge. If the best passage ends above the limit, or there is none, the search repeats at 20\,\%, or 50\,\% with a kite.

\begin{table}[tbp]
\centering
\footnotesize
\caption{Change of the saving, in points, against the finer or fuller search. All the search settings are the base except the one changed.}
\label{tab:routing_convergence}
\begin{tabular}{@{}p{8.5cm}p{4.5cm}@{}}
\toprule
Setting changed & Largest move \\
\midrule
Margin 3\,\%, time bin \SI{3}{h}, weather every \SI{6}{h} & 0.37 \\
Margin 10\,\%, final settings & 0.07 (\SI{140}{m} ship, rotors, Guadeloupe) \\
Margin 12\,\%, kite, coarser corridor & 0.11 (Kamsarmax, Guadeloupe) \\
Weather every \SI{3}{h} rather than \SI{6}{h} & 0.43 \\
Time bins of \SI{1.5}{h} rather than \SI{3}{h} & 0.63 \\
Every second speed against all 49, rotors and wings; the cost of a bare ship moves by at most 0.04\,\% & 0.06 \\
Every second speed, kite, coarser corridor, two of six savings & 0.31 \\
\bottomrule
\end{tabular}
\end{table}

The final pass stays within \SI{24}{h} of the path's arrival time at each stage. On 22 passages of the \SI{85}{m} ship each, in synthetic weather, every fourth speed missed the cost of all 49 by up to 0.97\,\% with wings and 2.7\,\% with the kite. Every second speed came within 0.03\,\%. A table of 25 speeds, or 9 commanded, understates the kite's saving by 0.26 and 0.34 points on average and by up to 1.3 points. Whole knots from \SI{4}{kn} understate the kite's saving by 0.39 points on average and by up to 0.73, the rotors' by 0.15 and the wings' by 0.06, and miss the cost of a bare ship by up to 0.57\,\%. 

\subsection{Thresholds of the lowest wind}
\label{app:min_wind_table}

For the base kites of mass $m$ and area $A$ the thresholds of \cref{sec:min_wind} are, in m/s:
\begin{center}
\begin{tabular}{rrrrrr}
\hline
$A$ [m$^2$] & $m$ [kg] & $V_\mathrm{cross}$ & $V_\mathrm{climb}$ & $V_\mathrm{lift}$ & launch \\
\hline
800 & 969 & 1.7 & 4.0 & 4.5 & 5.0 \\
1600 & 2172 & 1.8 & 4.2 & 4.5 & 5.0 \\
3200 & 5003 & 1.9 & 4.5 & 4.5 & 5.0 \\
\hline
\end{tabular}
\end{center}

\subsection{Tables, weather fields and routes}
\label{app:perf_grid}

The performance table has 49 speeds through the water spaced evenly from \SI{4}{kn} to the ship's maximum, continued at the same step down to \SI{1.5}{m/s}, true wind speed at \SI{10}{m} height in steps of \SI{1}{m/s} and true wind angle in steps of \SI{7.5}{\degree}. Kite forces are tabulated over apparent wind speeds from \SIrange{3}{30}{m/s} and courses every \SI{15}{\degree}. Weather fields are interpolated linearly in space and time over the Atlantic from \SI{35}{\degree S} to \SI{56}{\degree N} and from \SI{100}{\degree W} to \SI{10}{\degree E}. All routes start through the Channel on fixed waypoints. Corridor points reach \SI{2400}{nmi} to either side of the track, and the stages only require the ship to make headway towards its destination. Houston was Rotterdam's largest partner across the Atlantic by cargo weight at \SI{36}{Mt} in 2024, nearly all of it shipped to Rotterdam \citep{eurostat_mar_go_am_nl}.

\subsection{Forecast router}
\label{app:forecast_router}

The forecast router keeps a kite aloft where the plan has it and the real weather allows, since choosing the kite's state step by step ignores the launches that choice costs later: on a test passage that raised the cost by 0.28\,\%. Where the real weather leaves no route from the ship's position, or stops a step of its plan, the ship holds there on its hotel load with the kite stowed and re-plans on the next forecast, 6~hours later. After 72~hours at one point its passage fails. Over the 12 departures this happened once, for 6~hours, on a return passage of the rotor ship with its rotors out of service at the 3-day horizon, and every other cost is unchanged to four digits.

\subsection{Sourced values in the fuel choice and the method}
\label{app:fuel_pathways}

IRENA's 1.5\,\textdegree{}C pathway gives e-ammonia 4.5 times the shipping energy share of methanol in 2050 \citep{irena2021ship}. The Mærsk Mc-Kinney Møller Center's pathway grows ammonia from about 16\,\% in 2030 to just over half in 2050 \citep{mmmczcs2021its}. Rig masses per square metre are \SI{0.12}{t/m^2} for wings, \SI{0.63}{t/m^2} for rotors with their foundations and \SI{0.05}{t/m^2} for the kite with winch and mast. Rig areas span \SIrange{500}{16000}{m^2} for kites and a fifth to all of the deck capacity for wings and rotors. A kite's turn radius is \SI{45}{m} for \SI{125}{m^2}, scaled with the span. Deadweight is 72\,\% of displacement in the coaster family. Ship price grows with deadweight to the power 0.40, crew cost to the power 0.15, and the ammonia engine costs \SI{190}{USD/kW} more than a diesel. Gross tonnage follows the 1969 Tonnage Convention from the moulded volume up to a depth of 1.32 draughts.

\subsection{Details of the VTOL kite and the canopy in the method}
\label{app:method_extra}

The VTOL mass exponent of 2.75 lies between the 2.63 and 2.91 that Makani's small flown Wing~7 and the MX2 design imply, with bounds of 2.2 and 3.3 (\cref{eq:rigid_mass}). Its circles have a radius of at least 3.46 spans, and it flies only where its mean air speed pays for the climb to the top of the circle. Each of several kites gets a share of a lone kite's force from momentum theory, the theory of the disc a wind turbine sweeps (\cref{eq:loop_share}). Hovering costs electricity for \SI{6}{min} over each launch and landing. The cost models give USD~2.3 million for the \SI{54}{m^2} kite with its tether, winch and perch. Launch and landing together wear the tether as much as about \SI{8}{min} of flight. A canopy built for another force limit gets the cloth that keeps the same life, which grows with the limit only to the power 0.27, and seams and reinforcements in proportion to the limit.

\subsection{Launch mast height and mass}
\label{app:mast}

The mast height is 1.3 times the square root of the projected kite area, about 0.75 spans at the base aspect ratio; at a fixed aspect ratio the span is a constant times the square root of the area, so this rule is also one on span. The telescopic lift that raised the \SI{160}{m^2} kite of the Beluga SkySails to \SI{15}{m} \citep{wiwo2008}, read as \SI{125}{m^2} projected, gives 1.34. The \SI{35}{m} mast of the \SI{1000}{m^2} Airseas kite \citep{airseas_seawing_web}, whose area basis is not stated, gives 1.11 if projected and 1.25 if flat. Its steel mass is \SI{0.5}{t} plus a part that grows with kite area to the power $2/3$ times height to the power $5/3$.

\subsection{Prices off the production stage}
\label{app:price_stage}

Every source figure under the market factor is a cost before overhead and margin, and a bought-in part enters at its purchase price, which holds only its supplier's margin. The wing sail's base price covers its controls, hydraulics and assembly, and the prices of wings and rotors are twice their manufacturing cost. The composite is priced at a wind blade plant making about a thousand blades a year. A rotor lands within 13\,\% below and 6\,\% above Norsepower's 2019 prices of EUR~\num{300000} to \num{750000} per rotor \citep{tillig2020b}, or up to 22\,\% above them if those include installation, as their wording allows. The market factor on the tether would add 15\,\% to the wear per flight hour of a \SI{1600}{m^2} kite. The avionics of the VTOL kite, EUR~\num{150000} \citep{joshi2025}, are above the EUR~\num{15000} of a prototype, so the prototype figure, our low bound, is a kite flown without aviation certification and not an earlier stage. With every kite hardware and canopy price and the market factor at their upper bounds, close to the prices paid so far, the better of an \SI{800}{m^2} and a \SI{3200}{m^2} kite saves \tbd{6.6}\,\% on the \SI{140}{m} e-ammonia ship to Guadeloupe instead of \tbd{13.0}\,\% and \tbd{6.2}\,\% on the \SI{200}{m} ship instead of \tbd{13.5}\,\%, and on VLSFO it \tbd{costs 1 to 2\,\% more than it saves}.

\subsection{Kite losses and launch cost today and at the goal}
\label{app:launch_loss}

At today's estimate of 0.25\,\% per launch and landing, a canopy is lost about every 62 round trips on the \SI{140}{m} ship at 6.5 launches each, once in about seven years for one ship and about fourteen times a year in a fleet of a hundred. At the goal of 1 in \num{10000} the kite launches 8.4 times a round trip and the same fleet loses one in about a year and a half. One launch and landing of an \SI{800}{m^2} kite then costs USD~235 instead of USD~740, and of a \SI{3200}{m^2} kite USD~286 instead of USD~\num{2350}, of which USD~219 is the crew's time and the mean cost of turning the ship for a launch (\cref{sec:kite}).

\subsection{Figures of the published comparisons}
\label{app:published}

\begin{table}[tbp]
\centering
\caption{Figures of the studies in \cref{sec:related} that the text does not give.}
\label{tab:published_figures}
{\footnotesize
\begin{tabular}{@{}>{\raggedright\arraybackslash}p{3cm}>{\raggedright\arraybackslash}p{12.3cm}@{}}
\toprule
Study & Figures \\
\midrule
\citet{huang2025} & Rotors are 38\,\% of the rigs studied and 42\,\% of the installed systems. \\
\citet{traut2014} & \SI{500}{m^2} kite and one rotor on five routes: rotor 193 to \SI{373}{kW} on average, or 2 to 21\,\% of the propulsive power demand; kite 127 to \SI{461}{kW}, or 1 to 32\,\%. Kite lift coefficient 1.0 at a glide angle of \SI{15.9}{\degree}. \\
\citet{lu2020} & Aframax tanker of about \num{115000}~t on two Atlantic routes: fuel savings of 4.2 to 8.9\,\%. \\
\citet{rechevilanova2021} & A rotor gives 7.8 times the lift per square metre of a wing sail without flap. \\
\citet{podeur2016} & Kite lift coefficient 0.776 at a glide angle of \SI{12.2}{\degree}. \\
\citet{tillig2020} & Moving the rotors aft, so that the rudder carries more of the side force, raised the savings to 36 and 21\,\%. \\
\citet{tillig2020b} & Fixed speed of \SI{12}{kn}. \\
\citet{hansen2025} & Rotors cost at least \num{50000}~USD more per trip at a mean wind of \SI{5.4}{m/s} and saved up to about 10\,\% at \SI{10.7}{m/s}. \\
\citet{kramer2022} & Sail-induced resistance is 4.6 to 14.2\,\% of the total resistance. \\
\citet{erhard2013} & A \SI{320}{m^2} kite can replace up to \SI{2}{MW} of propulsion power. \\
\citet{skysails2018} & More than \num{1800} hours of air time over all its kites, from 6 to \SI{400}{m^2}. \\
\citet{bentin2016} & Multipurpose ship of \num{17500}~t deadweight between Baltimore and Wilhelmshaven, on the great circle and optimised routes. \\
\citet{mason2023} & \num{1080} simulated voyages on three routes, speed free between 8 and \SI{14}{kn}. Carbon saving 6.95--16.7\,\% on the great circle and 10.2--28.7\,\% if the departure forecast came true. On the eastbound leg of the longest route forecast errors removed 44\,\% of that extra saving; re-planning every 12 hours kept 79 to 88\,\% on the two long routes. \\
\citet{wendling2025} & The grid must not let the ship skip a stage within one time step; weather data at \SI{1}{\degree} or six hours, or coarser, understate the gain. \\
\citet{dhainaut2025} & Container ship of \num{2200}~TEU: three rotors per side saved 14\,\% of fuel, while 20 more hours of voyage time saved 30\,\% without rotors. \\
\citet{bentin2018} & Break-even fuel price, for an unspecified rig from another study: 500~EUR/t to pay back in 5 years or 250~EUR/t in 10. \\
\bottomrule
\end{tabular}}
\end{table}

\subsection{Published statements on kites and ship speed}
\label{app:speed_literature}

\citet{cedelft2016} find, at 100 and 85\,\% of service speed, that kites save relatively more at the lower speed and the same or less in absolute terms at the higher, and conclude that ``towing kites are best-suited for lower ship speeds'' (p.~48). The same report calls the need to slow down a barrier ``that has been overestimated so far'' for at least some systems (p.~10). \citet{traut2014} assume a 20\,\% speed reduction from the start and name slow and small ships as the first market. On a network of 15 routes they report savings of 1 to 21\,\% for a multi-purpose freighter at \SI{15}{kn} and 4 to 36\,\% at \SI{13}{kn}, from Schlaak et al. \citet{bentin2018} exclude ships of \SI{20}{kn} or more on grounds of wind state, for all rigs. \citet{emsa2023} state that kites do better at higher speeds, on one provider-reported pair of figures (p.~30), next to CE Delft's ranges.

Our kite saves under half of Podeur's figure at \SI{16}{kn}, the bottom of Bentin's great-circle range of 10 to 23\,\%, half or less of their routed 19 to 29\,\% and under a tenth of the 26.6\,\% of \citet{leloup2013}. Their kites pull with a glide ratio of about 4.6 to 4.7 \citep{podeur2016,leloup2013}, against our 3.9 for the whole system with tether and flight drag, which puts their window edge further upwind and their force higher across it. Podeur's ship moves only along its heading, in surge, and \citet{leloup2013} take the kite's drift force as negligible, so in both the side force costs nothing, while ours pays leeway, rudder drag and heel for it. The 26.6\,\% of \citet{leloup2013} is one eastbound crossing from New York to the Lizard at \SI{15.5}{kn} in one recorded set of winds, and Podeur's eastbound leg saves 18.2\,\% against 6.0\,\% westbound. Bentin's routed figures let the route grow by up to a fifth to save energy and take the kite's pull from an empirical law in the apparent wind \citep{bentin2018}.

